\section{Ethical and Regulatory Analysis}
\label{sec:ethics}

A grid-world simulation processes no personal data, harms no one, and needs
no licence. The analysis of this section is therefore explicitly
\emph{prospective}: it asks what would follow --- ethically and legally ---
if the capability this project prototypes were embodied in real aircraft
over real terrain. The exercise is not hypothetical hand-wringing. The task
specification of Section~\ref{sec:problem} --- find a hidden target in
minimum time, under partial observability, with a team resilient to
attrition --- is the abstract core of application profiles that range from
the unambiguously beneficial to the gravest imaginable, and the technical
analysis of Section~\ref{sec:mas} (in particular the autonomy placement of
Section~\ref{subsec:mas-autonomy} and the opacity trade-off of
Section~\ref{subsec:mas-beyond}) supplies exactly the vocabulary the
ethical questions are posed in. We proceed from the dual-use diagnosis
(Section~\ref{subsec:eth-dualuse}), through the two European regulatory
frameworks that would govern civilian deployments --- data protection
(Section~\ref{subsec:eth-gdpr}) and the AI Act
(Section~\ref{subsec:eth-aiact}) --- to the military scenario that escapes
both (Section~\ref{subsec:eth-laws}), the responsibility problem that
underlies all of it (Section~\ref{subsec:eth-responsibility}), and a set of
concrete design commitments (Section~\ref{subsec:eth-recommendations}).

\subsection{A dual-use capability}
\label{subsec:eth-dualuse}

Nothing in this system knows what its target \emph{is}. The environment
marks one cell; the policy learns to reach it quickly. Instantiate the
target as a missing hiker and the system is a search-and-rescue asset
\citep{queralta2020collaborative}; as a wildfire front, an environmental
monitor \citep{julian2019distributed}; as a person of interest, a
surveillance tool; as a military objective whose localization triggers
engagement, a component of a weapon system. The capability is
deployment-agnostic by construction --- and that is the precise sense in
which it is \emph{dual-use}: the same artifact serves civilian and military
ends without modification, the classification hinging entirely on context
of use.

Two features of this particular design sharpen the point beyond the generic
observation that search generalizes. First, the project's central technical
contribution --- fault tolerance through redundancy, decentralized fault
knowledge, and graceful degradation (Section~\ref{subsec:mas-resilience})
--- is exactly the property that makes a hostile swarm hard to defeat:
resilience to random attrition and resilience to defensive fire are the
same competence under different descriptions. Second, the single-policy
generalization across team sizes, sensing ranges, and communication regimes
(Section~\ref{subsec:train-dr}) is, read adversarially, robustness of the
threat across countermeasure conditions. One cannot celebrate these
properties in Section~\ref{sec:mas} and disown them here; the designer of a
capability owns its profile of uses, and the honest response is not denial
but analysis --- which uses are governed by what rules, where the rules run
out, and what the designer can do unilaterally
(Section~\ref{subsec:eth-recommendations}).

\subsection{Data protection: the GDPR lens}
\label{subsec:eth-gdpr}

Consider first the benign deployment: drones searching for missing persons
or assessing disaster zones over inhabited areas. Real perception is not an
oracle bit but cameras and sensors over public and private space, and
aerial imagery of identifiable individuals is personal data; its processing
falls squarely under the General Data Protection Regulation
\citep{gdpr2016} --- or, when the operator is a law-enforcement authority
acting in that capacity, under its sibling instrument for police
processing, Directive (EU) 2016/680. Four consequences of the GDPR's
architecture bear directly on a drone-search deployment.

\emph{Lawful basis and purpose limitation.} Consent is structurally
unavailable --- one cannot consent whoever happens to be under the flight
path --- so processing must rest on another Article~6 basis: protection of
vital interests fits the rescue scenario, public task fits civil
protection. The decisive discipline is purpose limitation: data collected
to find a missing person may not drift into general surveillance of
whoever else was observed --- \emph{mission creep} is not an abstract worry
for a system whose fused maps accumulate, by design, everything any team
member has ever seen (Section~\ref{subsec:prob-comm}).

\emph{Minimisation --- and an argument this architecture supplies.} The
GDPR requires that processing be limited to what the purpose necessitates.
Here the system's own abstraction makes a concrete privacy argument: every
mechanism of Sections~\ref{sec:problem}--\ref{sec:architecture} ---
coordination, fusion, dispersion, fault handling --- consumes only
\emph{occupancy-level} information (free, obstacle, target, crash site).
Nothing in the coordination layer needs imagery, identity, or appearance.
A privacy-respecting embodiment would therefore process raw sensor data
on-board, transiently, solely to update the occupancy and detection maps,
and share \emph{only those maps} between drones and with operators ---
data-protection-by-design in the sense of Article~25, implemented as an
architectural boundary between a perception layer (privacy-critical,
on-board, ephemeral) and the coordination layer this report describes
(abstract, shareable). The gap between what the swarm \emph{could} record
and what its task \emph{needs} is exactly the space the minimisation
principle demands be closed.

\emph{Impact assessment and accountability.} Systematic monitoring of
publicly accessible areas on a large scale triggers a mandatory data
protection impact assessment (Article~35(3)(c)) --- a multi-drone sweep of
an inhabited search area is a textbook instance. And the decentralized
architecture raises a non-trivial accountability question: the deploying
organization is plainly the controller, but the processing itself is
replicated across autonomous nodes whose knowledge states synchronize
opportunistically (Section~\ref{subsec:prob-comm}) --- every drone is a
processing location, retention must be enforced on every replica, and a
data-subject request must reach state that was fused, anonymously, into
multiple maps. The anonymity of fusion --- a coordination virtue in
Section~\ref{subsec:mas-coordination} --- becomes a governance burden the
moment the fused content is personal data: provenance-free knowledge
cannot be selectively erased.

\subsection{The EU AI Act}
\label{subsec:eth-aiact}

The AI Act \citep{aiact2024} governs AI systems by risk tier: prohibited
practices, high-risk systems subject to ex-ante requirements, and
lighter transparency regimes below. A learned multi-drone search policy is
unambiguously an AI system in the Act's sense; \emph{which} tier it
occupies depends, once again, on deployment.

\emph{Where civilian deployments land.} Used by or on behalf of public
authorities for locating persons in a law-enforcement context, the system
falls within the high-risk categories of Annex~III (law enforcement;
similarly migration and border management); deployed as a safety component
of critical infrastructure inspection, high-risk again. If the search
target is a \emph{person} and localization proceeds by identifying
individuals from sensor data, the system approaches remote biometric
identification --- whose real-time use in publicly accessible spaces for
law enforcement is prohibited outright (Article~5), save for narrowly
enumerated exceptions which, notably, include the \emph{targeted search
for missing persons}: the Act's drafters anticipated precisely the benign
reading of this capability and carved it out, subject to authorization and
safeguards. A civilian search system of this kind should therefore be
engineered, from the start, to the high-risk compliance envelope.

\emph{What the high-risk obligations would demand.} The obligations of
Chapter~III map onto this project's pipeline with almost uncomfortable
precision, and the mapping is worth stating because it shows the distance
between research code and deployable system. Risk management (Article~9)
would formalize what Section~\ref{subsec:train-correctness} does
informally. Data governance (Article~10) --- training data must be
relevant and representative of the operating domain --- is, for a policy
trained entirely in simulation, a question about the domain-randomization
ranges of Section~\ref{subsec:train-dr}: DR coverage \emph{is} the data
governance story, and its limits (Section~\ref{subsec:exp-discussion})
are compliance gaps. Record-keeping (Article~12) demands runtime logging;
the project's seeded reproducibility --- one seed regenerates layout,
faults, and hence the full episode
(Section~\ref{subsec:train-infra}) --- is a strong starting point that
would need tamper-evident deployment logging on top. Human oversight
(Article~14) is the sharpest gap: the system as built has \emph{no
runtime human role whatsoever} (Section~\ref{subsec:mas-autonomy}), so an
oversight interface --- monitoring, intervention, override --- would have
to be added, not merely documented. Accuracy and robustness (Article~15)
is the article this project can speak to most directly: the
fault-robustness and generalization evaluations of
Section~\ref{sec:experiments} are, in substance, Article~15 evidence ---
while their statistical character (performance over seeded samples, not
guarantees) marks the boundary of what such evidence can establish
(Section~\ref{subsec:eth-responsibility}).

\emph{The military carve-out.} All of the above evaporates in the scenario
the next subsection considers: Article~2(3) excludes from the Act's scope
systems placed on the market or used \emph{exclusively for military,
defence or national security purposes}. The European Union's flagship AI
regulation, including its prohibited-practices tier, simply does not apply
to a targeting swarm. The asymmetry deserves emphasis: the \emph{rescue}
deployment of this capability faces DPIAs, conformity assessments, and
oversight requirements; the \emph{lethal} deployment of the same
capability faces, from EU AI law, nothing. What fills the gap is the older
and thinner fabric of international humanitarian law and the slow-moving
intergovernmental process around autonomous weapons --- the subject of the
next subsection.

\subsection{The autonomous-weapons scenario}
\label{subsec:eth-laws}

Suppose, then, the ugly instantiation: the target is a military objective,
and detection cues engagement --- by a human downstream, or by the system
itself. In the latter case the assembly meets the standard definition of
an autonomous weapon system: one that, once activated, selects and engages
targets without further human intervention \citep{dod2012autonomy}. No
dedicated treaty governs such systems; the governance landscape consists
of international humanitarian law (IHL), the guiding principles affirmed
by the CCW Group of Governmental Experts --- IHL applies fully to
autonomous weapons, and human responsibility for their use cannot be
transferred to the machine \citep{ccw2019guiding} --- and normative
positions such as the ICRC's, which calls for ruling out unpredictable
autonomous weapons and those targeting persons, and for strict
constraints (geographic, temporal, situational, and of human supervision)
on the rest \citep{icrc2021position,sharkey2019autonomous}. Against that
landscape, the system described in this report would be an unacceptable
weapon component --- and the reasons are not generic misgivings about
``killer robots'' but specific, mechanistic properties established in the
preceding sections.

\emph{Distinction fails at the architecture's core.} IHL's first demand is
that attacks distinguish military objectives from civilians. The system's
target representation is a binary, monotone memory bit: once any drone
marks a cell as target, the mark persists, propagates through fusion to
the whole connected team, and --- because fusion is an anonymous maximum
with no provenance and no retraction
(Section~\ref{subsec:prob-comm}) --- \emph{cannot be corrected}. A single
false detection becomes, irreversibly, the entire team's truth; the very
monotonicity that makes cooperative mapping conflict-free makes
misidentification unrecallable. Add the belief staleness quantified in
Section~\ref{subsec:mas-resilience} --- drones act for whole intervals on
outdated pictures of the world --- and the mismatch with a legal standard
that requires case-by-case verification \emph{at the moment of attack} is
structural, not incidental.

\emph{Proportionality has no substrate.} The proportionality rule requires
weighing anticipated military advantage against expected incidental harm
--- a contextual, value-laden judgement. The policy's entire evaluative
apparatus is a Q-function distilled from the reward of
Table~\ref{tab:reward}: time, coverage, collision. There is no input
channel through which civilian presence could even \emph{enter} the
decision, let alone be weighed. This is the goal-oriented/goal-governed
distinction of Section~\ref{subsec:mas-agents} bearing legal weight: a
system whose objectives are distilled into weights cannot deliberate about
competing considerations, because there is nothing in it that represents
considerations at all.

\emph{Precaution collides with emergence.} IHL requires those who plan
attacks to take constant care and to cancel or suspend when circumstances
change. But the team-level behaviour of this system is, by the analysis of
Section~\ref{subsec:mas-emergence}, emergent: not specified anywhere,
produced by the interaction of learned policies, shaped but not dictated
by training. Meaningful human control --- in the influential analysis, the
requirement that the system \emph{track} human moral reasons and that its
conduct be \emph{traceable} to humans who understand it
\citep{santonidesio2018meaningful} --- fails on both prongs: tracking,
because the only ``reasons'' the system responds to are distilled search
economics; tracing, because no human can inspect \emph{why} the collective
did what it did (Section~\ref{subsec:eth-responsibility}). And the
operational profile at issue --- communication-denied environments,
exactly where autonomy is militarily attractive --- is the profile in
which a human \emph{could not} supervise even in principle: the same
range-limited communication that makes centralized control unimplementable
(Section~\ref{subsec:mas-why}) makes real-time human oversight
unimplementable too. The autonomy that solves the engineering problem is
the autonomy that dissolves the control requirement.

\emph{Generalization is not verification.} Finally, the system's
robustness claims are statistical: performance across domain-randomized,
seeded conditions (Section~\ref{sec:experiments}). An adversary is
precisely an out-of-distribution generator --- jamming pushes the
communication regime below anything trained, decoys exploit the
irreversible target memory, spoofed beacons corrupt the belief layer this
system trusts implicitly. Deploying against an adaptive opponent a policy
whose guarantees are distributional is unpredictability by construction
--- the property the ICRC position singles out as grounds for
prohibition \citep{icrc2021position}.

\subsection{Responsibility and opacity in self-organising systems}
\label{subsec:eth-responsibility}

Beneath both regulatory analyses lies one structural problem that this
report's own technical sections make unusually crisp: \emph{where, in a
system like this, is the seat of responsibility?} The classical conditions
for responsibility --- knowledge and control --- are strained twice over.

They are strained \emph{vertically} by learning. The behaviour ultimately
executed was produced by an optimization process; the designers chose the
reward, the ranges, and the architecture, but nobody chose --- or can
enumerate --- the policy's dispositions. This is the responsibility gap in
its canonical form \citep{matthias2004responsibility}: the manufacturer's
control ends where the learning begins, yet no one downstream gains the
control the manufacturer lost. The system's opacity is not incidental
polish that better documentation could remove: as
Section~\ref{subsec:mas-beyond} argued, a Q-network answers ``why'' with a
number, not a trace --- explainability was traded away, knowingly, for
adaptivity, and with it went the evidential substrate that after-the-fact
accountability (an investigation, an audit, a court) would need.

They are strained \emph{horizontally} by multi-agency. Outcomes here are
produced by four interacting policies, an environment's mediation, and the
accidents of who fused with whom when
(Section~\ref{subsec:mas-coordination}); the emergent regularities of
Section~\ref{subsec:mas-emergence} belong to no drone. Moral philosophy
knows this as the problem of many hands; this system exhibits its purest
technical form, since the ``hands'' are homogeneous, anonymous, and
individually thoughtless. When an emergent behaviour causes harm, the
candidate bearers of responsibility --- operator, commander, integrator,
policy trainer, reward designer --- each stand at a defensible distance
from the outcome, and the CCW principle that responsibility cannot pass to
the machine \citep{ccw2019guiding} names the constraint without resolving
the allocation.

Product-liability law patches part of the civil-damages corner of this
problem --- the revised EU directive extends strict liability for
defective products explicitly to software and AI systems, and to defects
arising from post-deployment learning \citep{pld2024} --- but liability
for compensation is the \emph{floor}, not the substance, of
accountability. What remains, and what
Section~\ref{subsec:eth-recommendations} operationalizes, is a design
obligation: if responsibility is to be attributable at all, the system
must be built so that human decisions remain the tracable causes of its
consequential actions --- responsibility-by-design as the governance
counterpart of privacy-by-design.

\subsection{Design recommendations}
\label{subsec:eth-recommendations}

The analysis above is not merely critical; it implies a concrete design
posture, most of it implementable unilaterally by the developer. Stated as
commitments for any real embodiment of this work:

\begin{itemize}[leftmargin=1.5em]
  \item \textbf{Detection, never engagement.} The system's output boundary
    is a \emph{location report to a human}; no actuator capable of harm is
    ever placed downstream of the policy. This single architectural
    commitment removes the system from the autonomous-weapon category by
    construction and preserves the human judgement that
    Section~\ref{subsec:eth-laws} showed the policy cannot supply.
  \item \textbf{A human on the loop, with real authority.} For any
    consequential downstream action, a human confirms on the basis of
    \emph{raw evidence} (imagery), not the system's binary target bit ---
    directly countering the irreversible-belief failure mode --- with
    abort authority whose availability is treated as a mission constraint:
    where communication denial would sever oversight, the mission profile,
    not the oversight, is what yields (cf.\ Article~14 of the AI Act).
  \item \textbf{Hard mission envelopes.} Geofencing, mission duration
    bounds, and target-class restrictions enforced \emph{outside} the
    learned component --- as environment-level constraints in the sense of
    Section~\ref{subsec:mas-environment}, whose compliance is verifiable
    by construction precisely because they are artefacts, not policies
    (Section~\ref{subsec:mas-artefacts}). This is the ICRC's
    constraint-based prescription \citep{icrc2021position} translated into
    this system's architecture.
  \item \textbf{Retraction-capable knowledge.} Replace the monotone target
    memory with a revisable representation (confidence-weighted, decaying,
    provenance-carrying), so that a false detection can be corrected and
    an erasure request honoured --- fixing in one stroke the
    distinction-critical failure of Section~\ref{subsec:eth-laws} and the
    erasure problem of Section~\ref{subsec:eth-gdpr}. Notably, this is a
    coordination-medium upgrade of exactly the kind
    Section~\ref{subsec:mas-beyond} recommended on engineering grounds.
  \item \textbf{Privacy-preserving perception.} On-board, ephemeral
    processing of raw sensing; only occupancy-level abstractions are
    stored, fused, or transmitted (Section~\ref{subsec:eth-gdpr}); imagery
    retention only upon human-confirmed detection, under the deployment's
    lawful basis.
  \item \textbf{Tamper-evident accountability logging.} Extend the seeded
    reproducibility of Section~\ref{subsec:train-infra} into deployment:
    signed logs of observations, beliefs, fusion events, and actions,
    sufficient to reconstruct \emph{what each drone knew and when} --- the
    minimal evidential basis Section~\ref{subsec:eth-responsibility} found
    missing.
  \item \textbf{Staged autonomy with demonstrated competence.} Deployment
    graduates through supervised levels (teleoperated, supervised
    autonomous, autonomous-within-envelope) in the manner of the
    established autonomy-level frameworks
    \citep{sae2021j3016,dod2012autonomy}, with promotion gated on the
    quantitative evidence of Section~\ref{sec:experiments} --- and on its
    honest limits.
\end{itemize}

None of this dissolves the dual-use character diagnosed in
Section~\ref{subsec:eth-dualuse} --- no design measure can prevent a
determined actor from rebuilding the capability without the safeguards.
What the measures achieve is narrower and still worth having: they make
the \emph{artifact as shipped} lawful in its intended deployments,
auditable when it fails, and useless as a weapon without substantial,
deliberate, and therefore attributable modification.
